\section*{第 \ref{ch:g4:animal-reproduction} 章——动物如何繁殖}

\begin{solution}{exo:g4:animal-reproduction:1}
一个\omterm{def:g4:animal-reproduction:sexual}{精子}细胞与一个\omterm{def:g4:animal-reproduction:sexual}{卵细胞}结合，由此开始
形成一个新的\omterm{def:g1:animals-around-us:animal}{动物}幼体。\omterm{def:g4:animal-reproduction:sexual}{卵细胞}来自\omterm{def:g4:animal-reproduction:sexual}{雌性}，\omterm{def:g4:animal-reproduction:sexual}{精子}
细胞来自\omterm{def:g4:animal-reproduction:sexual}{雄性}。
\end{solution}

\begin{solution}{exo:g4:animal-reproduction:2}
\omterm{def:g4:animal-reproduction:ovivivi}{卵生}：\omterm{def:g4:animal-reproduction:sexual}{雌性}产卵，幼体在卵中发育完成，
在母体之外——母鸡、青蛙（还有鳟鱼、
蜥蜴、蜗牛、\omterm{prop:g3:sorting-living-things:groups}{昆虫}）。\omterm{def:g4:animal-reproduction:ovivivi}{胎生}：幼体在母体内生长，
由母体供养，活着出生——母马、
鲸（几乎所有\omterm{prop:g3:sorting-living-things:groups}{哺乳动物}都是如此）。
\end{solution}

\begin{solution}{exo:g4:animal-reproduction:3}
\omterm{def:g4:animal-reproduction:sexual}{雄性}鸣叫，把\omterm{def:g4:animal-reproduction:sexual}{雌性}召到池塘边；每对
将各自的细胞排放到水中，\omterm{def:g4:animal-reproduction:sexual}{精子}在水中与卵相遇——
即\omterm{def:g4:animal-reproduction:sexual}{受精}；裹在胶质中的\omterm{def:g4:animal-reproduction:sexual}{受精}卵发育成了
六月的\omterm{ex:g2:animals-grow-up:frog}{蝌蚪}。
\end{solution}

\begin{solution}{exo:g4:animal-reproduction:4}
鳟鱼的\omterm{def:g4:animal-reproduction:sexual}{受精}发生在水中、雌鱼体外——卵
与\omterm{def:g4:animal-reproduction:sexual}{精子}一同排出。母鸡的\omterm{def:g4:animal-reproduction:sexual}{受精}发生在体内，
在\omterm{prop:g4:animal-reproduction:where}{交配}时进行。在陆地上，裸露的细胞会干燥；在水中，
它们可以在体外安全相遇。
\end{solution}

\begin{solution}{exo:g4:animal-reproduction:5}
小鸡在有壳的卵中、母鸡体外生长，靠
卵中贮存的卵黄为食。小马驹在母马体内
生长十一个月，由母体持续供养，活着出生。
\end{solution}

\begin{solution}{exo:g4:animal-reproduction:6}
一个物种产生的幼体越多，每个幼体得到的照料就越少；
越少，得到的照料就越多。两个极端：鳕鱼
的数百万枚无人照料的卵，以及大象
仅有的一头小象受守护长达十年。
\end{solution}

\begin{solution}{exo:g4:animal-reproduction:7}
因为成功只要求平均而言，有足够多的
幼体存活下来、足以替代亲代。在数百万枚卵中，总有
少数能躲过海洋中的种种危险——庞大的数量
\emph{正是}保护。
\end{solution}

\begin{solution}{exo:g4:animal-reproduction:8}
幼体数量少（比起数千只，五只已算少），在
母体内发育，由母体哺育和守护：猫明显属于
少而受照料的一类——与\omterm{prop:g3:sorting-living-things:groups}{哺乳动物}总体一样。
\end{solution}

\begin{solution}{exo:g4:animal-reproduction:9}
它从\emph{每一}个亲代那里都获得了某些
东西——幼体并不只是其母亲的复制品。双亲
都共同决定了小马驹的样子。
\end{solution}

\begin{solution}{exo:g4:animal-reproduction:10}
产卵量少——只有几十或几百枚卵，而非
数百万枚。照料与数量此消彼长：守护、
扇风和保卫的父亲，能让一小窝卵度过
那些危险，而无人照料的卵需要数百万枚才能做到。
\end{solution}

\begin{solution}{exo:g4:animal-reproduction:11}
\omterm{def:g4:animal-reproduction:ovivivi}{卵生}——且无人照料：埋下的一窝卵只能
靠沙子的温度孵化，刚孵出的幼体独自
冲向大海（数量庞大正是它们的策略）。她必须在陆地上产卵，
因为\omterm{prop:g4:classifying-animals:classes}{爬行动物}的卵有壳且需呼吸空气——若产在
海里，卵就会被淹死。决定这一规则的，是她所属纲的特征，
而非她的习性。
\end{solution}
